\documentclass[11pt,a4paper]{article}
\usepackage[T1]{fontenc}
\usepackage{lmodern,amsmath,amssymb,amsthm,mathtools,microtype,booktabs}
\usepackage[margin=27mm]{geometry}
\usepackage[colorlinks=true,linkcolor=blue,citecolor=blue,urlcolor=blue]{hyperref}
\setlength{\emergencystretch}{2em}
\newtheorem{theorem}{Theorem}[section]
\newtheorem{lemma}[theorem]{Lemma}
\newtheorem{proposition}[theorem]{Proposition}
\newcommand{\Q}{\mathbb Q}\newcommand{\Z}{\mathbb Z}
\DeclareMathOperator{\Tr}{Tr}
\title{Even abscissa valuations and a smaller eleven-variable degree bound}
\author{}\date{17 November 2026}
\begin{document}\maketitle
\begin{abstract}
The first four tests of the derivative detector have even valuation
at every consumed odd place, including for unrestricted auxiliary
elliptic points. They therefore need no additional squaring before
balancing. This preserves extraction, fixed-place and contact preparation,
and all finite-character guards. A complete eleven-variable ordinary-$abc$
family has degree at most $230932097$, improving the previous bound
$232094337$. Its seven-contact variant has degree at most $356146817$
and a degree-seven physical lifting equation. For a halfpoint fixture,
the actual weight equation has an explicit genus-twelve model with
smaller coefficients. Exact physical compositions have genera $233$
and $181$ for nine and seven contacts, respectively. Global actual-trace
realization, and hence completeness in nine variables, remains unproved.
\end{abstract}

\section{An even-valuation lemma}
The relevant assertion is local parity, not a global square class.
\begin{lemma}\label{lem:general}
Let $L$ be a number field, let $v$ be a normalized nonarchimedean
valuation, and let
\[
                   E:y^2=x(x^2+A x+B)
\]
with $v(A)\ge0$ and $v(B)=0$. Every finite point $(x,y)\in E(L)$
with $xy\ne0$ satisfies $v(x)\in2\Z$.
\end{lemma}
\begin{proof}
If $v(x)>0$, then $B$ is the unique term of least valuation in
$x^2+A x+B$, so $2v(y)=v(x)$. If $v(x)<0$, then $x^2$ is the
unique term of least valuation, giving $2v(y)=3v(x)$. Each equality
forces $v(x)$ to be even. The case $v(x)=0$ is immediate.
\end{proof}

For the retained source
\[
 K=\Q(\sqrt2),\qquad E:y^2=x^3-dx,\qquad
 d=11-6\sqrt2,\qquad N_{K/\Q}(d)=49,
\]
the lemma applies at every place above a rational prime other than
seven. All consumed odd primes are outside the existing exceptional
set $S_0$, which contains two and seven; they are unramified in $K$.
Thus their normalized slot valuations of rational $u$ equal $v_p(u)$.
The same parity conclusion applies to every guarded difference point.

\begin{lemma}\label{lem:balanced}
Let $p\notin S_0$, put $\nu=v_p(u)$, and suppose $\nu$ is odd.
If $C\in K^\times$ has even valuation in both slots, then
\[
                 D=\frac{C}{1+uC},\qquad
                  v(D)=\max\{v(C),-\nu\}
\]
at each place above $p$. In particular, if $\nu>0$, then
$v(D)>0$ if and only if $v(C)>0$. If $\nu<0$, then
$v(D^2/u)\ge-3\nu>0$.
\end{lemma}
\begin{proof}
The integer $\nu+v(C)$ is odd and therefore nonzero, so the two
terms of $1+uC$ have unequal valuations. Hence
$v(D)=v(C)-\min(0,\nu+v(C))$, which is the displayed maximum.
The remaining assertions follow directly. No rational or global
$K$-square assumption on $C$ is used.
\end{proof}

\section{The modified source}
Retain ordinary rational $abc$, all arithmetic premises of the derivative
construction \cite{Derivative}, and its full auxiliary extension
\cite{FullAux}. In particular,
\[
 2E(K)=\Z P,\qquad
 P=\left(\frac{57-28\sqrt2}{8},\frac{-660+417\sqrt2}{32}\right).
\]
The root $A=[n]P$ retains its two-coordinate chart and integral index.
The two auxiliaries $H,U$ may range over all guarded affine points.
With $f=1/x$ and $\zeta=-x/y$, put
\[
 r=\frac{\zeta(U)}{\zeta(H)},\qquad
 r'=\overline{\frac{\zeta(U-A)}{\zeta(H-P)}},
\]
\[
 (C_1,C_2,C_3,C_4,C_5,C_6)
 =\bigl(f(H),f(U),\overline{f(H-P)},\overline{f(U-A)},r-r',r-X\bigr).
\]
All original point, ordinate, difference and quotient guards remain.
By Lemma~\ref{lem:general}, the first four tests have even valuations
at every consumed place. The last two tests need not have this property
and may vanish.

For the degree-improving variant use
\begin{equation}\label{eq:height}
 \beta=X^2+B,\quad T=\Tr x(A),\quad M=\beta^4+T^4,
 \quad a=M^4,\quad t=M^8-z^2,
\end{equation}
where $B\ge3$ is a fixed odd integer and $z\in\Q$ is the one free
slope. Use the nine shifts
$\mathcal S_9=\{-10,-8,-5,1,2,3,4,6,7\}$, and put
\[
 F=\prod_{c\in\mathcal S_9}((a+c)^2-z^2),\qquad
 u=-18F/F_a,\qquad w=\frac{u}{(1+u)^2}.
\]
The modification is exactly
\begin{equation}\label{eq:newD}
 D_i=\begin{cases}
 C_i/(1+uC_i),&1\le i\le4,\\
 C_i^2/(1+uC_i^2),&5\le i\le6.
 \end{cases}
\end{equation}
The detector and intrinsic input remain
\begin{equation}\label{eq:detector}
 W=\prod_{i=1}^6(1+D_i^2/u),\quad V=\Tr(W^{-1}),\quad
 b=4w/V^2,\quad\eta=7b.
\end{equation}
Retain all original nonzero requirements on $M,t,F,F_a,u,1+u$,
the detector product and trace, and $\eta$. The first four balanced
norm guards are replaced by the norms of $1+uC_i$; the last two are
unchanged. We do not omit any guard on the ground that its value might
usually be nonzero.

Two free pole coordinates $v,s$ and the finite root of
\begin{equation}\label{eq:pole}
 Y^2=\eta(\eta+v^2)^2s^4-\eta(\eta-v^2)^2(1-\eta)^2
\end{equation}
complete the eleven-coordinate source. The final root may vanish in
soundness; all preceding point radicands and syntactic denominators
remain units.

\section{Arithmetic transfer}
\begin{theorem}\label{thm:transfer}
The source \eqref{eq:height}--\eqref{eq:pole} is sound and complete
for every sufficiently large fixed odd $B$, under the same retained
premises as \cite{FullAux}. It therefore defines $\Z$ by a polynomial
in eleven total variables, including the input. The nine-variable
actual-trace source obtained by replacing \eqref{eq:pole} with
$Y^2=V^2+w(2-V)$ remains sound; its completeness is not asserted.
\end{theorem}
\begin{proof}
Let $\nu=v_p(u)>0$ be odd, $p\notin S_0$, and let $q\le0$
be the largest slot valuation of $W$. Every factor of $W$ has
valuation $\min(0,2v(D_i)-\nu)$, since its two terms have distinct
valuations when the second is nonzero. The intrinsic pole implication
excludes negative odd valuations of $b$, giving
$1\le v_p(b)\le\nu+2q$. Thus in a slot attaining $q$, each factor
has valuation greater than $-\nu/2$, and $v(D_i)>\nu/4>0$.
Lemma~\ref{lem:balanced} forces $v(C_i)>0$ for $i\le4$.
For the comparison tests the old identity
$v(D_i)=\max(2v(C_i),-\nu)$ gives the same conclusion. The retained
formal-comparison theorem therefore extracts $X\equiv n\pmod p$
and $p$-integrality of $X$ exactly as before.

Neither the contact weight $u$ nor its height argument has changed.
The complete height absorption argument of \cite{Derivative} therefore
applies, with the same positive margin $44272779/198801800$.
The extension to half-integral auxiliary indices is exactly that of
\cite{FullAux}; only the four local abscissa parities were additionally
used. They hold on the full auxiliary domain, not merely on old witnesses.

For completeness retain the old multiplication subfamily
$U=[n]H$, $H\in2E(K)$. At a selected contact the original first four
tests have valuation at least two. By Lemma~\ref{lem:balanced}, their
new $D_i$ still have valuation at least two when $v_p(u)=1$.
Hence their factors are one modulo $p$, and the original successful-slot
and unit-trace proof is unchanged. Comparison factors are unchanged.
At a negative odd valuation of $u$, Lemma~\ref{lem:balanced} makes
each of the first four factors one modulo $p$ regardless of abscissa
poles. The old squared balancing does the same for the last two.
Thus $V\equiv2$, and the proof that $\eta$ has square reduced
denominator remains valid, including new denominator primes of $F_a$.
At even valuations the parity argument for $b$ is unchanged.

At each fixed place, make $z$ large relative to the already prepared
nonzero tests. Both versions of \eqref{eq:newD} then satisfy
$D_i\sim1/u$ for nonzero $C_i$, while identically zero comparisons
remain zero. Thus $W\to1$, $V\to2$, and the original two-adic and
seven-adic preparation remains valid. The finite norm guards also stay
nonzero after taking the original neighborhoods sufficiently small and
the slope sufficiently large.

For completeness, all prepared point characters can still be chosen
real with positive first four tests. Indeed all old integral-index
points have square nonzero abscissae. The algebraic identity
\[
 x(Q)x(R)x(Q-R)
   =\left(\frac{x(Q)y(R)+x(R)y(Q)}{x(Q)-x(R)}\right)^2
\]
shows that their required differences have positive abscissae at each
real sign character whenever the original difference guards hold.
This identity remains valid for either ordinate sign. Therefore
$1+uC_i>0$ for $i\le4$ when $u>0$, and the last two balanced
denominators are also positive. Every detector factor is at least one,
so $0<V\le2$. The original real preparation gives $0<\eta<1$.
This argument is for the subfamily used in completeness; arbitrary
full-auxiliary real points are not assumed to have positive abscissae.

The intrinsic-pole existence and all-character finite-norm avoidance
proofs now apply unchanged, since their inputs have exactly the same
preparation. No preceding radicand is allowed to vanish. The final
root remains permitted to vanish in soundness.

Finally the actual-trace implication excluding negative odd valuations
of $b$ depends only on $w=u/(1+u)^2$ and $V$, not on the formula for
the balanced tests. It therefore supplies the same soundness for the
nine-variable source. At every prepared contact, negative odd valuation,
fixed place and real character, its local-square proof transfers as
well. Those local statements do not prove global rational squareness.
\end{proof}

\section{Integral specification and degree}
The direct identities
\[
 \frac{1/x}{1+u/x}=\frac1{x+u},\qquad
 \frac{\Delta^2/X_0}{1+u\Delta^2/X_0}
                  =\frac{\Delta^2}{X_0+u\Delta^2}
\]
avoid unnecessary fraction expansions. Here a difference abscissa is
stored as $X_0/\Delta^2$. The original nonzero guards on $x$, $X_0$
and $\Delta$ make these identities and the corresponding norm-guard
replacements equivalent at every finite-algebra character.
The coefficient circuit uses these identities directly for the first
four factors; its independent native checker evaluates
$C_i/(1+uC_i)$ instead.

The three finite point blocks have rank $64$; adjoining the terminal
quadratic root gives rank $128$. The unchanged guarded subalgebra norms
and terminal-root separator \cite{Terminal} use one encoder coordinate.
The coordinate count is $1+6+1+2+1=11$. The source has 1865 arithmetic
nodes and 32 native guards. Independent specification traversal gives
\[
 \deg F_B\le230932097<2^{28},\qquad
 \#\operatorname{supp}(F_B)<2^{264},\qquad
                     \operatorname{content}(F_B)=1.
\]
This is an upper bound for the finite determinant specification, not
an expanded exact degree. The coefficient bound is computable for every
specified $B$. For the stored fixture $B=3$, its coefficient $\ell^1$
norm is at most $2^{480197504}$. No sufficient numerical defining $B$
is claimed.

The same modification applies to the seven-contact source of
\cite{Seven}, where $a=M^8$, $t=M^{16}-z^2$ and
$\mathcal S_7=\{-7,-6,-4,1,3,5,8\}$. Its strict height margin is
$38695561/597203000$. It has 1815 nodes, 32 native guards, rank 128,
and degree at most $356146817<2^{29}$, with fewer than $2^{270}$
monomials. Its purpose is the degree-seven physical lift, not a lower
final degree. Each variant is checked at 384 finite-field sign characters,
and independent corrupted-degree and missing-guard controls are rejected.

\section{The new actual halfpoint curve}
Use the guarded full-auxiliary fixture
\[
 X=1,\quad B=3,\quad A=P,\quad H=U=R_0=(1+\sqrt2,2\sqrt2).
\]
One has $2R_0=P$ and $H-P=-H$. The actual six tests are
$(c,c,\bar c,\bar c,0,0)$ with $c=\sqrt2-1$. In this section the
first four factors use the unsquared balancing \eqref{eq:newD}.
Define
\begin{align*}
 D_0(u)&=u^2(1-2u-u^2)^2,\\
 R_0(u)&=D_0(u)+2u^3+4u^2+6u+1,\\
 \Pi_0(u)&=4D_0(u)^4(1+u)^2
       +2uR_0(u)^2\bigl(R_0(u)^2-D_0(u)^2\bigr).
\end{align*}
The point $R_0$ and polynomial $R_0(u)$ are distinguished by their
arguments. Since $c+\bar c=-2$ and $c\bar c=-1$,
\[
 J_0(u)=\frac{\bigl(u(1+uc)^2+c^2\bigr)
                  \bigl(u(1+u\bar c)^2+\bar c^2\bigr)}
                  {u^2(1-2u-u^2)^2}
       =\frac{R_0(u)}{D_0(u)}.
\]
It follows exactly that $W=J_0^2$, $V=2D_0^2/R_0^2$, and
\begin{equation}\label{eq:curve}
 Y_0^2=\Pi_0(u),\qquad Y_0=(1+u)R_0(u)^2Y
\end{equation}
is the actual weight cover.

\begin{proposition}\label{prop:genus}
The curve \eqref{eq:curve} has genus twelve after normalization.
Every rational weight $u\ne0,-1$ satisfies the selected-character
balanced, detector-product and trace denominator guards. For the seven-
and nine-contact physical maps, the squared-slope covers have genera
$90$ and $116$, respectively; the slope covers have genera $181$ and
$233$, respectively.
\end{proposition}
\begin{proof}
The degree of $\Pi_0$ is 26, its leading coefficient is four, and
zero is a simple root. The stored finite-field Bezout certificate
proves it squarefree. Thus its genus is twelve.
The quadratic $1-2u-u^2$ has nonsquare discriminant eight. The monic
polynomial $R_0$ has constant coefficient one and
$R_0(1)=17$, $R_0(-1)=1$; hence it has no rational root. Consequently
$D_0R_0\ne0$ at every rational $u\ne0$. The denominator and trace
formulas prove the assertion about the selected character.

Let $k$ be seven or nine. The fixed physical scale is
$a=(10621537/256)^8$ when $k=7$ and
$a=(10621537/256)^4$ when $k=9$. Differentiate
$F(a,T_0)=\prod_{c\in\mathcal S_k}((a+c)^2-T_0)$ before specializing
$a$, and form the common-content primitive pair
$A_k/B_k=-2kF/F_a$ over $\Z[T_0]$. Its degrees are $k,k-1$.
The polynomial
\[
              P_k(T_0)=B_k(T_0)^{26}\Pi_0(A_k(T_0)/B_k(T_0))
\]
has degree $26k$, is squarefree, and does not vanish at zero; the
certificate reconstructs the pair and verifies the appropriate
finite-field Bezout identity independently. Thus
$Y_1^2=P_k(T_0)$ has genus $13k-1$,
$Y_2^2=T_0P_k(T_0)$ has genus $13k$, and
$Y_1^2=P_k(z^2)$ has genus $26k-1$.
\end{proof}

A rational point $(u,Y_0)$ on \eqref{eq:curve} still needs a root of
$uB_k(T_0)-A_k(T_0)$ in $\Q$, with $B_k(T_0)\ne0$ and
$T_0$ a rational square. The physical guard $a^2-T_0\ne0$ and
all original point, whole-character and preparation requirements remain.
The known boundary points $u=0,\infty$ do not meet these conditions.
No admissible rational point is asserted here. The fixture $B=3$ is
not certified to exceed a sufficient defining threshold. In particular
none of these arithmetic reductions proves nine- or twenty-one-variable
completeness.

\section{Reproduction}
The directory
\nolinkurl{computation/few_variables/even_valuation11_v75}
contains both complete integer circuit specifications, their pinned
producers, and independently adapted degree and native point checks.
The second directory
\nolinkurl{computation/pole_parameters/even_halfpoint_trace_v75}
contains the exact weight polynomials, both physical pairs, and their
finite-field certificates. Its independent reconstruction imports no
producer routine and differentiates before specializing the physical
scale. These finite computations establish their listed algebraic
assertions; the mathematical transfer proofs above retain the source's
stated infinite arithmetic premises.

\begin{thebibliography}{9}
\bibitem{Derivative}\emph{An eleven-variable construction under rational $abc$},
\nolinkurl{papers/conditional_few_variables/derivative_detector11_20261117_v63.tex}.
\bibitem{FullAux}\emph{Full auxiliary elliptic points in the eleven-variable source},
\nolinkurl{papers/conditional_few_variables/full_auxiliary11_20261117_v71.tex}.
\bibitem{Seven}\emph{Seven-contact detectors and a genus-twelve actual weight curve},
\nolinkurl{papers/conditional_few_variables/seven_contact_detector_20261117_v74.tex}.
\bibitem{Terminal}\emph{A terminal quadratic root and complete polynomial specifications},
\nolinkurl{papers/finite_algebra/terminal_quadratic_root_20261115_v40.tex}.
\end{thebibliography}
\end{document}
